\subsection{Emergence of the Schr\"{o}dinger equation as an effective projected dynamics}
  \label{app:schrodinger-derivation}

  \subsubsection*{Main theorem}

    The following theorem states the complete result.
    Steps~1--5 and Appendix~\ref{app:schrodinger-hamiltonian} constitute its proof.

    \begin{theorem}[Emergence of quantum dynamics from projected substrate relaxation]
      \label{thm:schrodinger-emergence}
      Let $(\Omega, \mathcal{R}, \Pi)$ be a Cosmochrony triple consisting of:
      \begin{itemize}
        \item[(H1)] a space $\Omega$ of admissible relational configurations with a nonlinear
        relaxation operator $\mathcal{R}$,
        \item[(H2)] a generally non-injective projection $\Pi : \Omega \to O$ with finite
        projective resolution $\varepsilon_\Pi > 0$,
        \item[(H3)] a Born--Infeld admissibility bound: every mode satisfies
        $A_n(\tau) \leq c_\chi/\sqrt{\lambda_n}$ at all $\tau$,
        \item[(H4)] a quasi-stable background $\bar{\chi}$ with $\|\delta\chi\| \ll 1$
        and $A_n \ll A_n^{\max}$,
        \item[(H5)] a spectrally admissible projected sector: the low-lying modes of
        $\delta\chi$ diagonalize $\mathcal{L} \equiv D\mathcal{R}|_{\bar\chi}$.
      \end{itemize}
      Then:
      \begin{enumerate}
        \item[(C1)] \textbf{Anti-Hermitian reduction.}
        The restriction of $\mathcal{L}$ to the projected-stable sector is anti-Hermitian:
        \[
          \mathcal{L}\big|_{\mathrm{proj}} = -i\,\Omega,
          \qquad \Omega = \Omega^\dagger,\quad \sigma(\Omega) \subset \mathbb{R}_{\geq 0}.
        \]

        \item[(C2)] \textbf{Unique modal evolution.}
        For each mode $u_n$, the unique solution of the projected dynamics is
        $c_n(\tau) = c_n(0)\,e^{-i\omega_n\tau}$,
        giving the modal Schr\"{o}dinger equation
        \[
          i\hbar_{\mathrm{eff}}\,\partial_\tau c_n = E_n\,c_n,
          \qquad E_n = \hbar_{\mathrm{eff}}\,\omega_n.
        \]

        \item[(C3)] \textbf{Effective Schr\"{o}dinger equation.}
        The projected wavefunction $\psi(x,t) = \sum_n c_n(t)\,\phi_n(x)$ satisfies
        \[
          i\hbar_{\mathrm{eff}}\,\partial_t\psi
          = \hat{H}_{\mathrm{eff}}\,\psi + \mathcal{N}_{\Pi,\mathrm{BI}}[\psi],
        \]
        where $\hat{H}_{\mathrm{eff}}$ is Hermitian, defined spectrally by
        $\hat{H}_{\mathrm{eff}}\phi_n = E_n\phi_n$,
        and $\mathcal{N}_{\Pi,\mathrm{BI}} = O(A_n/A_n^{\max},\,\ell_\chi^2\Delta_g)$.

        \item[(C4)] \textbf{Standard Hamiltonian in the non-relativistic limit.}
        In the additional regime $\ell_\chi^2\Delta_g \ll 1$ and locally homogeneous substrate,
        \[
          \hat{H}_{\mathrm{eff}}
          = -\frac{\hbar_{\mathrm{eff}}^2}{2m_{\mathrm{eff}}}\,\Delta_g
          + V_{\mathrm{eff}}(x)
          + O(\ell_\chi^2\,\Delta_g^2),
        \]
        where $\Delta_g$ is the Laplace--Beltrami operator of the emergent metric $g^{\mu\nu}$
        defined by $\sigma_2(-\Delta_\chi)(x,k) = g^{\mu\nu}(x)\,k_\mu k_\nu$,
        and $m_{\mathrm{eff}}$ is determined by the spectral gap of the mode family.

        \item[(C5)] \textbf{Structural status of unitarity.}
        Norm preservation on the projected-stable sector, and hence unitarity of the
        effective dynamics, is not an independent postulate.
        It holds if and only if (H2) and (H3) hold simultaneously:
        \[
          \textup{unitarity}
          \;\Longleftrightarrow\;
          (\textup{Born--Infeld capacity}) \cap (\textup{finite projective resolution}).
        \]
      \end{enumerate}
    \end{theorem}

    \begin{proof}
      (C1) is established in Step~3: $\operatorname{Re}(\mu_n) > 0$ is excluded by (H3);
      $\operatorname{Re}(\mu_n) < 0$ is excluded by (H2) since the mode amplitude falls below
      $\varepsilon_\Pi$ at finite $\tau$; hence $\operatorname{Re}(\mu_n) = 0$ is the only
      structurally admissible value, giving $\mu_n = -i\omega_n$ and
      $\mathcal{L}|_{\mathrm{proj}} = -i\Omega$.

      (C2) follows by direct integration of $\partial_\tau c_n = -i\omega_n c_n$, which is a
      first-order linear ODE with constant coefficients whose unique global solution is
      $c_n(\tau) = c_n(0)\,e^{-i\omega_n\tau}$.
      Multiplication by $\hbar_{\mathrm{eff}}$ and the energy-frequency relation
      $E_n = \hbar_{\mathrm{eff}}\omega_n$
      (Section~\ref{subsec:energy-frequency-solitons}) give
      $i\hbar_{\mathrm{eff}}\,\partial_\tau c_n = E_n\,c_n$.

      (C3) follows from (C2) by reconstruction: setting
      $\psi = \sum_n c_n\phi_n$ and applying $i\hbar_{\mathrm{eff}}\partial_t$,
      \[
        i\hbar_{\mathrm{eff}}\,\partial_t\psi
        = \sum_n E_n\,c_n(t)\,\phi_n(x) = \hat{H}_{\mathrm{eff}}\psi
      \]
      at leading order.
      The correction $\mathcal{N}_{\Pi,\mathrm{BI}}$ collects Born--Infeld nonlinearities
      $O(A_n/A_n^{\max})$ from (H3) and inter-modal coupling from (H4); both are
      negligible under (H4).

      (C4) follows from the spectral identification in
      Appendix~B.\ref{app:schrodinger-hamiltonian}:
      $\omega_n^2 = \lambda_n$ at leading order (from the admissibility
      envelope~\cite{Beau2026a1}), combined with the operator--metric correspondence
      $\sigma_2(-\Delta_\chi)(x,k) = g^{\mu\nu}k_\mu k_\nu$~\cite{Beau2026c}, and the
      non-relativistic expansion
      $\\hbar_{\mathrm{eff}}\sqrt{-\Delta_g}
      \simeq m_{\mathrm{eff}}c^2 + \frac{-\hbar_{\mathrm{eff}}^2\Delta_g}{2m_{\mathrm{eff}}}$.

      (C5) is the logical converse of (C1): $\operatorname{Re}(\mu_n) = 0$ requires both
      (H2) and (H3); if either fails, modes with $\operatorname{Re}(\mu_n) \neq 0$ enter
      the sector, breaking norm preservation.
    \end{proof}

    \begin{remark}
      Hypothesis~(H5) --- spectral admissibility of the projected sector --- is a structural
      assumption on the relational substrate, not an additional quantization postulate.
      It asserts that the low-lying modes of $\mathcal{L}$ are diagonalizable within the
      projected-stable sector.
      This is guaranteed whenever the substrate background $\bar{\chi}$ is spectrally
      admissible in the sense of~\cite{Beau2026a1}.
    \end{remark}

    \begin{remark}
      The theorem makes explicit that two distinct regimes are involved.
      Conclusions~(C1)--(C3) require only~(H1)--(H5) and hold at the level of general
      projected substrate dynamics.
      Conclusion~(C4) requires the additional continuum-metric hypothesis implicit in
      Step~B of Appendix~B.\ref{app:schrodinger-hamiltonian}.
      The non-relativistic Hamiltonian $-\hbar^2\Delta_g/(2m)$ is therefore a further
      specialization, not the most general output of the theorem.
    \end{remark}

    We now develop the theorem step by step.
    The purpose of the following construction is not to introduce additional assumptions,
    but to make explicit how the effective Schr\"odinger dynamics follows from projected
    substrate relaxation in the weak-amplitude regime.

    Let $\Omega$ denote the space of admissible relational configurations and
    \[
      \Pi : \Omega \to O
    \]
    the generally non-injective projection defining the effective observables.
    A projected state does not correspond to a single microscopic configuration, but to an
    equivalence class of configurations under $\Pi$.
    The effective state therefore encodes only the resolvable modal content that survives the
    projection at finite projective resolution.

    \paragraph{Step 1: Substrate relaxation and linearization.}

      At the substrate level, the primitive dynamics takes the generic relaxation form
      \[
        \partial_{\tau}\chi = \mathcal{R}[\chi],
      \]
      where $\tau$ denotes the intrinsic ordering parameter of the substrate and $\mathcal{R}$ is
      the nonlinear relaxation operator.
      In the weak-amplitude regime around a quasi-stable background configuration $\bar{\chi}$,
      we write $\chi = \bar{\chi} + \delta\chi$ with $\|\delta\chi\| \ll 1$.
      Linearizing gives
      \[
        \partial_{\tau}\delta\chi = \mathcal{L}\,\delta\chi + O(\delta\chi^2),
        \qquad
        \mathcal{L} \equiv D\mathcal{R}\big|_{\bar{\chi}}.
      \]

    \paragraph{Step 2: Modal decomposition in the spectrally admissible sector.}

      We assume that the projected regime is spectrally admissible, so that the relevant sector of
      $\delta\chi$ can be expanded on a discrete family of effective modes:
      \[
        \delta\chi(\tau) = \sum_n c_n(\tau)\,u_n,
        \qquad
        \mathcal{L}\,u_n = \mu_n\,u_n.
      \]
      In the stable dissipative sector of the full substrate, $\operatorname{Re}(\mu_n) \leq 0$.
      When $\mu_n$ is real and negative, one recovers a purely damped mode
      $c_n(\tau) \propto e^{\mu_n \tau}$.
      This is the generic substrate behavior and does not by itself produce quantum-like dynamics.

    \paragraph{Step 3: Restriction to the projected-stable sector and anti-Hermitian reduction.}

      The key step is the restriction to the \emph{projected-stable sector}: the subspace of modes
      that contribute a finite and persistently resolvable effective amplitude to the observable
      sector $O$ at finite projective resolution.
      Membership in this sector is governed by two independent structural constraints inherited
      from the Cosmochrony framework.

      The first constraint is the Born--Infeld admissibility bound established in the spectral
      admissibility programme~\cite{Beau2026a1}.
      A mode with Laplacian eigenvalue $\lambda_n$ is dynamically admissible if and only if its
      effective amplitude satisfies
      \[
        A_n(\tau) \leq A_n^{\max} = \frac{c_\chi}{\sqrt{\lambda_n}}
      \]
      at all values of $\tau$.
      In the linearized regime, $A_n(\tau) = A_n(0)\,e^{\operatorname{Re}(\mu_n)\,\tau}$.
      If $\operatorname{Re}(\mu_n) > 0$, the amplitude grows without bound in finite intrinsic time,
      violating the Born--Infeld bound.
      Growing modes are therefore strictly excluded from the admissible sector.

      The second constraint is projectability.
      The projection $\Pi$ identifies configurations within the same equivalence class
      $\Pi^{-1}(o)$, and possesses a finite projective resolution: two configurations are
      indistinguishable under $\Pi$ whenever their effective amplitudes differ by less than the
      resolvability threshold of the projection.
      A mode with $\operatorname{Re}(\mu_n) < 0$ has a decaying effective amplitude and therefore
      becomes indistinguishable from the trivial configuration once it falls below the resolvability
      threshold of $\Pi$, contributing nothing distinguishable to any effective observable at the
      given projective resolution.
      This exclusion does not rely on the asymptotic limit $\tau \to \infty$.
      It holds at every finite value of $\tau$ at which the amplitude has decayed below the
      projective resolution scale.
      Decaying modes therefore do not belong to the projected-stable sector.

      The two constraints together leave $\operatorname{Re}(\mu_n) = 0$ as the only structurally
      admissible possibility for modes in the projected-stable sector.
      The eigenvalues therefore take the form $\mu_n = -i\omega_n$ with $\omega_n \in \mathbb{R}$.
      Equivalently, the restriction of $\mathcal{L}$ to the projected-stable sector is anti-Hermitian:
      \[
        \mathcal{L}\big|_{\mathrm{proj}} = -i\,\Omega,
      \]
      where $\Omega$ is a Hermitian operator with real spectrum $\{\omega_n\}$.

    \begin{remark}
      The two constraints play complementary and non-redundant roles.
      The Born--Infeld bound eliminates growing modes through a dynamical capacity argument:
      the substrate cannot sustain unbounded flux.
      Projectability eliminates decaying modes through a resolvability argument: a mode whose
      effective amplitude falls below the projective resolution threshold is not a degree of freedom
      of the effective description at finite projective resolution.
      Neither constraint alone is sufficient; their conjunction is what leaves
      $\operatorname{Re}(\mu_n) = 0$ as the only structurally admissible value.
      Norm preservation in the projected-stable sector is therefore not an additional postulate.
      It is the necessary consequence of bounded substrate capacity and finite projective resolution
      acting simultaneously.
      Equivalently, this establishes that unitarity of the effective quantum dynamics is not a
      fundamental symmetry postulate, but the structural signature of the intersection of two
      independent constraints:
      \[
        \textup{unitarity}
        = (\textup{Born--Infeld capacity}) \cap (\textup{non-injective projection}).
      \]
      \end{remark}

    \paragraph{Step 4: From anti-Hermitian operator to phase evolution ---
    the closure argument.}

      The anti-Hermitian reduction $\mathcal{L}|_{\mathrm{proj}} = -i\,\Omega$ established
      in Step~3 determines the modal evolution completely.
      For any mode $u_n$ in the projected-stable sector, the linearized equation reads
      \[
        \partial_\tau c_n = -i\,\omega_n\,c_n,
        \qquad \omega_n \in \mathbb{R},\; \omega_n \geq 0.
      \]
      This is a first-order linear ordinary differential equation with constant coefficients,
      whose unique global solution is
      \[
        c_n(\tau) = c_n(0)\,e^{-i\omega_n\tau}.
      \]
      The decomposition $c_n(\tau) = a_n(\tau)\,e^{-i\omega_n\tau}$ with slowly varying
      $a_n$ is therefore not an ansatz but a rewriting of the exact solution, with $a_n$
      absorbing only residual corrections from nonlinear and inter-modal terms:
      \[
        \partial_\tau a_n
        = \bigl[\partial_\tau c_n + i\omega_n c_n\bigr] e^{i\omega_n\tau}
        = O\!\left(a_n^2,\,\nabla a_n,\,\Pi_{\mathrm{sat}}\right).
      \]
      In the weak-amplitude, approximately injective regime these correction terms are
      negligible, so $a_n$ is constant to leading order.

      Multiplying the exact modal equation by $\hbar_{\mathrm{eff}}$ and using
      $E_n = \hbar_{\mathrm{eff}}\,\omega_n$, one obtains
      \[
        i\hbar_{\mathrm{eff}}\,\partial_\tau c_n = E_n\,c_n.
      \]
      This is the modal Schr\"{o}dinger equation.
      It follows from $\mathcal{L}|_{\mathrm{proj}} = -i\,\Omega$ by direct integration,
      with no additional hypothesis.

    \paragraph{Step 5: Reconstruction of the wavefunction and the operator equation.}

      Passing from the intrinsic ordering parameter $\tau$ to the effective projected time
      coordinate $t$ amounts, in the weak-curvature regime, to a monotone reparametrization
      $\partial_{\tau} = \zeta(x,t)\,\partial_t$ with $\zeta \simeq 1$ in a locally
      Minkowskian domain.
      Absorbing $\zeta$ into the local normalization of $\hbar_{\mathrm{eff}}$ gives, to
      leading order,
      \[
        i\hbar_{\mathrm{eff}}\,\partial_t c_n = E_n\,c_n.
      \]
      Let $\phi_n(x)$ denote the effective projected spatial eigenmodes.
      The projected state reconstructs as
      \[
        \psi(x,t) = \sum_n c_n(t)\,\phi_n(x).
      \]
      Applying $i\hbar_{\mathrm{eff}}\,\partial_t$ to both sides gives
      \[
        i\hbar_{\mathrm{eff}}\,\partial_t\psi
        = \sum_n \bigl(i\hbar_{\mathrm{eff}}\,\partial_t c_n\bigr)\phi_n(x)
        = \sum_n E_n\,c_n(t)\,\phi_n(x)
        = \hat{H}_{\mathrm{eff}}\,\psi,
      \]
      where the effective Hamiltonian is defined spectrally by
      $\hat{H}_{\mathrm{eff}}\phi_n \equiv E_n\phi_n$.
      This yields
      \[
        \boxed{i\hbar_{\mathrm{eff}}\,\partial_t\psi = \hat{H}_{\mathrm{eff}}\,\psi.}
      \]

    \begin{proposition}[Effective Schr\"{o}dinger equation]
    Under hypotheses \textup{(H1)--(H5)} of Theorem~\ref{thm:schrodinger-emergence},
    the projected wavefunction $\psi(x,t) = \sum_n c_n(t)\,\phi_n(x)$ satisfies
    \[
      i\hbar_{\mathrm{eff}}\,\partial_t\psi
      = \hat{H}_{\mathrm{eff}}\,\psi + \mathcal{N}_{\Pi,\mathrm{BI}}[\psi],
    \]
    where $\hat{H}_{\mathrm{eff}}$ is Hermitian and spectrally defined by
    $\hat{H}_{\mathrm{eff}}\phi_n = E_n\phi_n$, and
    $\mathcal{N}_{\Pi,\mathrm{BI}} = O(A_n/A_n^{\max},\,\ell_\chi^2\Delta_g)$.
    In the regime $A_n \ll A_n^{\max}$ and $\ell_\chi^2\Delta_g \ll 1$, the correction
    vanishes and the equation reduces to the standard linear Schr\"{o}dinger equation.
      \end{proposition}

    \begin{proof}
      Conclusion~(C2) of Theorem~\ref{thm:schrodinger-emergence} gives
      $i\hbar_{\mathrm{eff}}\,\partial_t c_n = E_n\,c_n$.
      Reconstruction via $\psi = \sum_n c_n\phi_n$ and the spectral definition of
      $\hat{H}_{\mathrm{eff}}$ give~(C3).
      The correction $\mathcal{N}_{\Pi,\mathrm{BI}}$ is negligible under~(H4).
      \end{proof}

    \paragraph{Structural conclusions.}

      This derivation clarifies the status of the standard quantum formalism.

      First, the wavefunction $\psi$ is not fundamental.
      It is the compressed descriptor of a spectrally admissible family of projected substrate
      modes belonging to the projected-stable sector at finite projective resolution.

      Second, $\hbar_{\mathrm{eff}}$ is not a primitive parameter.
      It measures the finite projective resolvability of the effective description and appears
      because phase transport remains observable after non-injective compression.

      Third, the factor $i$ does not require a fundamental complex ontology.
      It is the algebraic consequence of the anti-Hermitian reduction of
      $\mathcal{L}|_{\mathrm{proj}}$, which itself follows from the joint action of the
      Born--Infeld bound and finite projective resolution.

      Fourth, linearity is not exact at all scales.
      It holds only when three conditions are simultaneously satisfied: the substrate remains
      sufficiently close to a quasi-stable background, Born--Infeld-type saturation effects
      remain negligible, and the projection remains approximately injective over the support
      of the state.
      Outside this regime,
      \[
        i\hbar_{\mathrm{eff}}\,\partial_t \psi
        = \hat{H}_{\mathrm{eff}}\,\psi + \mathcal{N}_{\Pi,\mathrm{BI}}[\psi],
      \]
      where $\mathcal{N}_{\Pi,\mathrm{BI}}$ collects the nonlinear contributions induced by
      bounded projective capacity and substrate saturation.
      The ordinary Schr\"{o}dinger equation is recovered as the universal infrared transport
      law of projected coherence within the projected-stable sector.

      The previous subsection establishes the emergence of a Schr\"odinger-type evolution
      equation at the operator level.
      The next subsection identifies the spectral origin of the corresponding Hamiltonian and
      derives its non-relativistic Laplace--Beltrami form.
